\begin{pastexam}[label={px:ncu-111-dsa-10}]{中央資工所 111 學年度 複選題第 10 題}
（複選；每題 5 分，全對才給分，答錯倒扣 1 分）

The adjacency list below is for an AOE network. The \textbf{end} field points to a list of adjacent
vertices, \textbf{dur} field is the duration of the activity, \textbf{link} field points to another
adjacent vertex, \textbf{vertex} field is the id of the adjacent vertex, \textbf{count} field is the
number of immediate predecessors. Which of the following statements are true?
\begin{enumerate}[label=(\Alph*)]
\item The total duration of the critical path is 17.
\item Path 0, 1, 4, 6, 8 is not a critical path.
\item Count fields from vertex 0 to vertex 8 should be 0, 1, 1, 1, 2, 2, 1, 2, 2.
\item The latest time that event 4 (i.e., vertex 4) can occur is 7.
\end{enumerate}

\begin{center}
\begin{tikzpicture}[x=1cm,y=1cm,>={Stealth[length=2mm]},font=\small]
\def\rh{0.56}% row pitch
\def\cw{0.6}% cell width
\def\ch{0.4}% list-node cell height
\node[font=\small\bfseries,anchor=base] at (0.6,0.55) {count end};
\node[font=\small\bfseries,anchor=base] at (3.05,0.55) {vertex dur link};
\foreach \i in {0,...,8} {
  \node[anchor=east] at (-0.2,-\i*\rh) {\i};
  \draw (0,-\i*\rh-\rh/2) rectangle (\cw,-\i*\rh+\rh/2);
  \draw (\cw,-\i*\rh-\rh/2) rectangle (2*\cw,-\i*\rh+\rh/2);
}
\node at (1.5*\cw,-8*\rh) {0};
\def\lnode#1#2#3#4#5{%
  \draw (#1,-#2*\rh-\ch/2) rectangle ++(3*\cw,\ch);
  \draw (#1+\cw,-#2*\rh-\ch/2) -- ++(0,\ch);
  \draw (#1+2*\cw,-#2*\rh-\ch/2) -- ++(0,\ch);
  \node at (#1+0.5*\cw,-#2*\rh) {#3};
  \node at (#1+1.5*\cw,-#2*\rh) {#4};
  \node at (#1+2.5*\cw,-#2*\rh) {#5};
}
\foreach \i in {0,...,7} \draw[->] (1.5*\cw,-\i*\rh) -- (2.05,-\i*\rh);
\lnode{2.1}{0}{1}{5}{}
\lnode{4.7}{0}{2}{4}{}
\lnode{7.3}{0}{3}{5}{0}
\draw[->] (2.1+2.5*\cw,0) -- (4.65,0);
\draw[->] (4.7+2.5*\cw,0) -- (7.25,0);
\lnode{2.1}{1}{4}{1}{0}
\lnode{2.1}{2}{4}{3}{0}
\lnode{2.1}{3}{5}{2}{0}
\lnode{2.1}{4}{6}{9}{}
\lnode{4.7}{4}{7}{7}{0}
\draw[->] (2.1+2.5*\cw,-4*\rh) -- (4.65,-4*\rh);
\lnode{2.1}{5}{7}{5}{0}
\lnode{2.1}{6}{8}{2}{0}
\lnode{2.1}{7}{8}{5}{0}
\end{tikzpicture}
\end{center}
\end{pastexam}

\begin{pastsoln}
答案是 (B)、(D)。

\textbf{讀出網路。} 第 $i$ 列的串列列出從事件 $i$ 出發的活動（頂點、持續時間），link 欄的 $0$ 表示串列
結束。 共 $11$ 個活動：
\[
\begin{array}{l}
0 \to 1\ (5),\quad 0 \to 2\ (4),\quad 0 \to 3\ (5),\quad 1 \to 4\ (1),\\
2 \to 4\ (3),\quad 3 \to 5\ (2),\quad 4 \to 6\ (9),\quad 4 \to 7\ (7),\\
5 \to 7\ (5),\quad 6 \to 8\ (2),\quad 7 \to 8\ (5).
\end{array}
\]
事件 $0$ 是起點，事件 $8$ 是終點（它的 end 欄是 $0$，沒有出邊）。

\textbf{最早時間 $ee$}（資結書第~22.5~節；依拓樸順序 $0, 1, \ldots, 8$，取所有前驅的最大值）：
\[
\arraycolsep=3pt
\begin{array}{lll}
ee[0] = 0, & ee[1] = 5, & ee[2] = 4, \\
ee[3] = 5, & ee[4] = \max(5 + 1,\; 4 + 3) = 7, & ee[5] = 5 + 2 = 7, \\
ee[6] = 7 + 9 = 16, & ee[7] = \max(7 + 7,\; 7 + 5) = 14, & ee[8] = \max(16 + 2,\; 14 + 5) = 19.
\end{array}
\]
\textbf{最晚時間 $le$}（從 $le[8] = ee[8] = 19$ 往回，取所有後繼的最小值）：
\[
\arraycolsep=3pt
\begin{array}{lll}
le[8] = 19, & le[7] = 19 - 5 = 14, & le[6] = 19 - 2 = 17, \\
le[5] = 14 - 5 = 9, & le[4] = \min(17 - 9,\; 14 - 7) = 7, & le[3] = 9 - 2 = 7, \\
le[2] = 7 - 3 = 4, & le[1] = 7 - 1 = 6, & le[0] = \min(6 - 5,\; 4 - 4,\; 7 - 5) = 0.
\end{array}
\]
活動 $i \to j$ 是關鍵活動當且僅當 $ee[i] = le[j] - \text{dur}$：$0 \to 2$、$2 \to 4$、$4 \to 7$、$7 \to 8$
四個，關鍵路徑是 $0, 2, 4, 7, 8$，長度 $4 + 3 + 7 + 5 = 19$。
\begin{itemize}
\item (A) 錯：關鍵路徑長 $19$，不是 $17$。
\item (B) 對：$0, 1, 4, 6, 8$ 的長度是 $5 + 1 + 9 + 2 = 17 < 19$，不是關鍵路徑（$1 \to 4$ 的
$ee[1] = 5 \ne le[4] - 1 = 6$，有一單位的寬裕）。
\item (C) 錯：count 是入分支度，依序是 $0, 1, 1, 1, 2, 1, 1, 2, 2$。 事件 $5$ 只有 $3 \to 5$ 一個前驅，
count 是 $1$，選項寫成 $2$。
\item (D) 對：$le[4] = 7$。
\end{itemize}
\end{pastsoln}
